\documentclass[aps,prc,onecolumn,superscriptaddress,nofootinbib,10pt]{revtex4-2}
\usepackage{amsmath,amssymb}
\usepackage{booktabs}
\usepackage[colorlinks=true,linkcolor=blue,citecolor=blue,urlcolor=blue]{hyperref}

\newcommand{\nB}{n_B}
\newcommand{\Msun}{M_\odot}
\newcommand{\rsun}{r_\odot}
\newcommand{\rhosol}{\rho_\odot}

\begin{document}

\title{Stellar structure in \texttt{eos/astro/tov}: \\
       the TOV system, tidal deformability, the crust, and rotation}

\author{M.~Guerrini}
\affiliation{University of Ferrara}

\begin{abstract}
\texttt{eos/astro/tov} turns an equation of state into a star. It integrates
the Tolman--Oppenheimer--Volkoff
system~\cite{Tolman1939,OppenheimerVolkoff1939} for the mass and radius, the
$l = 2$ static tidal perturbation~\cite{Hinderer2008,DamourNagar2009} for the
Love number $k_2$ and the dimensionless deformability $\Lambda$, and --- through
the external RNS code~\cite{StergioulasFriedman1995} --- uniformly rotating
models. It also assembles the table it integrates, joining a tabulated crust to
the core an \texttt{eos} model produced.

This note states every equation the code solves and every quantity it returns.
Two things in it are easy to get wrong and are given in full: the
\emph{dimensionless variables} the integration is actually carried out in,
because the residual expressions look unfamiliar until the scaling is written
down; and the \emph{jump} in the tidal variable at a first-order interface,
whose denominator is $m + 4\pi r^3 P$ and not $m$~\cite{TakatsyKovacs2020}.

This layer consumes tables and produces observables. It never imports model
internals, and the table it takes --- \texttt{EOSTable\_for\_TOV}, three parallel
arrays $(P, \varepsilon, \nB)$ --- lives in \texttt{eos/general/state.py},
which is the layer both a model and this one may import.
\end{abstract}

\maketitle

\section{What is integrated, and in which variables}
\label{sec:units}

The TOV system in geometric units ($G = c = 1$) is
%
\begin{align}
\frac{dm}{dr} &= 4\pi r^2 \varepsilon ,
\label{eq:dmdr_phys}\\[2pt]
\frac{dP}{dr} &= -\,\frac{(\varepsilon + P)\,(m + 4\pi r^3 P)}
                          {r\,(r - 2m)} .
\label{eq:dpdr_phys}
\end{align}
%
The code does not integrate these in physical units. It uses
%
\begin{equation}
\hat r = \frac{r}{\rsun}, \qquad
\hat m = \frac{m}{\Msun}, \qquad
\hat p = \frac{P}{P_c},
\label{eq:scaling}
\end{equation}
%
where $\rsun = 2G\Msun/c^2 = 2.953$~km is the solar Schwarzschild radius,
$\Msun = 1.11542\times10^{60}$~MeV$/c^2$, and $P_c$ is the \emph{central}
pressure of the star being built, so that $\hat p = 1$ at the centre. The
combination that appears throughout is
%
\begin{equation}
\rhosol \;=\; \frac{3\,\Msun}{4\pi\,\rsun^3} ,
\label{eq:rhosol}
\end{equation}
%
the mean density of a solar mass inside its own Schwarzschild radius. The
choice $\rsun = 2G\Msun/c^2$ is what makes the metric factor collapse:
%
\begin{equation}
1 - \frac{2Gm}{rc^2} \;=\; 1 - \frac{\hat m}{\hat r} ,
\label{eq:metric}
\end{equation}
%
with no constant left over.

\paragraph{One expression, two factors --- read this before comparing with the
source.} $4\pi r^2 \varepsilon$ appears both as a \emph{mass per unit radius}
(in $dm/dr$) and as a \emph{dimensionless} geometric quantity (inside $F$ and
$Q$ of Sec.~\ref{sec:tidal}), and it does NOT carry the same numerical factor
in the two places. A mass scales as
%
\begin{equation}
\frac{4\pi r^3 P}{\Msun} = 3\,\frac{P}{\rhosol}\,\hat r^3 ,
\qquad
\frac{d\hat m}{d\hat r} = 3\,\frac{\varepsilon}{\rhosol}\,\hat r^2 ,
\label{eq:scale_mass}
\end{equation}
%
directly from Eq.~\eqref{eq:rhosol}. A dimensionless geometric combination
picks up one further factor of $\rsun/\Msun = 2G/c^4$, which halves it:
%
\begin{equation}
\big[\,4\pi r^2 \varepsilon\,\big]_{\rm geom}
   \;=\; \frac{3}{2}\,\frac{\varepsilon}{\rhosol}\,\hat r^2 .
\label{eq:scale_geom}
\end{equation}
%
This is why the code carries $3/2$ in $F$ and in the first term of $r^2Q$ while
carrying $3$ in $d\hat m/d\hat r$ and in $m + 4\pi r^3 P$. Deriving one from
the other and expecting the same coefficient is the mistake to avoid.

Equations
\eqref{eq:dmdr_phys}--\eqref{eq:dpdr_phys} are therefore integrated as
%
\begin{align}
\frac{d\hat m}{d\hat r} &= 3\,\frac{\varepsilon}{\rhosol}\,\hat r^2 ,
\label{eq:dmdr}\\[2pt]
\frac{d\hat p}{d\hat r} &=
  -\,\frac{1}{2}\,\frac{\varepsilon}{P_c}\,\hat m\,
   \left(1 + \frac{P}{\varepsilon}\right)
   \left(1 + 3\,\frac{P_c}{\rhosol}\,\frac{\hat p\,\hat r^3}{\hat m}\right)
   \frac{1}{\hat r\,(\hat r - \hat m)} ,
\label{eq:dpdr}
\end{align}
%
with $P = \hat p P_c$ and $\varepsilon = \varepsilon(P)$ read from the table by
linear interpolation on the $P$ grid. Both right-hand sides are returned as
zero for $\hat p \le 0$ or $\hat r \le 0$, and $d\hat p/d\hat r$ is set to zero
if $|\hat r - \hat m| < 10^{-30}$, which is the horizon the integration must
never reach.

\paragraph{Initial conditions and termination.}
Integration starts at $\hat r_{\min} = 10^{-3}$ with the uniform-core value
%
\begin{equation}
\hat m(\hat r_{\min}) = \frac{4\pi}{3}\,
   \frac{\varepsilon_c\, (\hat r_{\min}\rsun)^3}{\Msun}
 \quad\text{written as}\quad
   \varepsilon_c \frac{4}{3}\pi \hat r_{\min}^3 / \Msun
\label{eq:m0}
\end{equation}
%
in the module's units, and $\hat p(\hat r_{\min}) = 1$. The surface is located
by a terminal event, $\hat p = 10^{-10}$, crossed from above; the solver is
\texttt{DOP853} with \texttt{rtol} $=10^{-10}$, \texttt{atol} $=10^{-12}$ and
dense output. The returned radius is $R = \hat r_s \rsun$ in km and the mass
$M = \hat m(\hat r_s)$ in $\Msun$, both already in the units
Eq.~\eqref{eq:scaling} defines.

\section{Baryonic mass}
\label{sec:mb}

The baryon number inside the star counts proper volume, not coordinate volume:
%
\begin{equation}
M_b \;=\; m_N \int_0^{R} \frac{4\pi r^2\, \nB(r)}
                              {\sqrt{1 - 2Gm(r)/rc^2}}\; dr
      \;=\; m_N \int_0^{\hat r_s} \frac{4\pi (\hat r \rsun)^2\, \nB}
                              {\sqrt{1 - \hat m/\hat r}}\; \rsun\, d\hat r ,
\label{eq:Mb}
\end{equation}
%
with $m_N$ the nucleon mass and $\nB(r) = \nB(P(r))$ from the table. The
integral is evaluated on 500 equally spaced points by the trapezoidal rule on
the dense output of the structure solve. The metric factor is floored at
$0.01$ where $\hat r \le \hat m$, which cannot happen at a converged solution
and guards the sampling only. $M_b - M$ is the gravitational binding energy.

\section{Tidal deformability}
\label{sec:tidal}

\subsection{The perturbation equation}

A static $l = 2$ even-parity perturbation of the metric obeys a second-order
equation for $H(r)$, which is integrated here in the logarithmic-derivative
variable
%
\begin{equation}
Y(r) \;=\; \frac{r\,H'(r)}{H(r)} ,
\label{eq:Ydef}
\end{equation}
%
because $Y$ stays finite where $H$ grows over many orders of magnitude. It
satisfies the Riccati form~\cite{Hinderer2008,DamourNagar2009,Postnikov2010}
%
\begin{equation}
\boxed{\;
r\,\frac{dY}{dr} + Y^2 + Y\,F(r) + r^2 Q(r) \;=\; 0 , \qquad Y(0) = 2 , \;}
\label{eq:Y}
\end{equation}
%
with
%
\begin{align}
F(r) &= \frac{1 - 4\pi r^2\,(\varepsilon - P)}{1 - 2m/r} ,
\label{eq:F}\\[2pt]
r^2 Q(r) &= \frac{4\pi r^2\left[\,5\varepsilon + 9P
             + \dfrac{\varepsilon + P}{dP/d\varepsilon}\,\right]}{1 - 2m/r}
          \;-\; \frac{6}{1 - 2m/r}
          \;-\; \left[\frac{m + 4\pi r^3 P}{r\,(1 - 2m/r)}\right]^2 .
\label{eq:Q}
\end{align}
%
$Y(0) = 2$ is the regular solution at the centre. In the variables of
Sec.~\ref{sec:units}, writing $\hat e = \varepsilon/\rhosol$ and
$\pi_P = P/\varepsilon$, the code evaluates
%
\begin{align}
F &= \Big(1 - \tfrac{3}{2}\,\hat e\,\hat r^2\,(1 - \pi_P)\Big)\,
     \frac{\hat r}{\hat r - \hat m} ,
\label{eq:Fcode}\\[2pt]
r^2 Q &= \tfrac{3}{2}\,\hat e\,\hat r^2\,\frac{\hat r}{\hat r - \hat m}
         \left[\,5 + 9\pi_P + \frac{1 + \pi_P}{c_s^2}\,\right]
       \;-\; \frac{6\,\hat r}{\hat r - \hat m}
       \;-\; \frac{\big(\hat m + 3(P_c/\rhosol)\,\hat r^3 \hat p\big)^2}
                   {(\hat r - \hat m)^2} ,
\label{eq:Qcode}
\end{align}
%
and steps $d Y/d\hat r = -(Y^2 + Y F + r^2 Q)/\hat r$ alongside
Eqs.~\eqref{eq:dmdr}--\eqref{eq:dpdr}, so the three components
$(\hat m, \hat p, Y)$ are advanced by one solver.

\paragraph{The sound speed is tabulated, not differentiated in the loop.}
$c_s^2 = dP/d\varepsilon$ enters Eq.~\eqref{eq:Q} through
$(\varepsilon+P)/(dP/d\varepsilon)$. It is computed ONCE before the
integration, as a finite difference of $\varepsilon(P)$ on the table's own $P$
grid, $c_s^2 = 1/(d\varepsilon/dP)$, clipped to $[10^{-6}, 1]$, and then
interpolated like any other column. Differentiating inside the right-hand side
would make every step depend on a numerical derivative of an interpolant. The
lower clip matters at a Maxwell plateau, where $d\varepsilon/dP \to \infty$ and
$c_s^2 \to 0$ exactly.

\subsection{The Love number}

With the surface value $Y_R = Y(R)$ and the compactness $C = GM/Rc^2$,
%
\begin{align}
k_2 &= \frac{8}{5}\,C^5\,\frac{(1-2C)^2\,\big[\,2 - Y_R + 2C(Y_R-1)\,\big]}{D} ,
\label{eq:k2}\\[4pt]
D &= 2C\Big[\,6 - 3Y_R + 3C(5Y_R - 8)\,\Big] \nonumber\\
  &\quad + 4C^3\Big[\,13 - 11 Y_R + C(3Y_R - 2) + 2C^2(1 + Y_R)\,\Big]
     \nonumber\\
  &\quad + 3(1-2C)^2\Big[\,2 - Y_R + 2C(Y_R-1)\,\Big]\ln(1-2C) ,
\label{eq:k2den}\\[4pt]
\Lambda &= \frac{2}{3}\,\frac{k_2}{C^5} .
\label{eq:Lambda}
\end{align}
%
This is Eq.~(23) of Ref.~\cite{Hinderer2008}. The implementation returns
$k_2 = \Lambda = 0$ outside $0.005 < C < 0.5$ and where $|D| < 10^{-30}$: the
closed form is ill-conditioned there, and a number produced from it would not
mean anything.

\subsection{The jump at a first-order interface}
\label{sec:jump}

Where the energy density is discontinuous while $P$ is continuous --- the
surface of a self-bound star, or an internal Maxwell transition --- $H'$ is
discontinuous and $Y$ jumps. Integrating $Y$ straight through such an
interface is simply wrong. The jump is
%
\begin{equation}
\boxed{\;
\Delta Y \;=\; -\,\frac{4\pi r_d^3\, \Delta\varepsilon}
                       {m(r_d) + 4\pi r_d^3 P(r_d)}
 \;=\; -\,\frac{3(\Delta\varepsilon/\rhosol)\,\hat r_d^3}
               {\hat m + 3 (P_d/\rhosol)\,\hat r_d^3} \;}
\label{eq:dY}
\end{equation}
%
with $\Delta\varepsilon = \varepsilon_{\rm in} - \varepsilon_{\rm out} > 0$
across the interface at $r_d$.

\textbf{The denominator is $m + 4\pi r^3 P$, not $m$.} The two agree only at a
surface, where $P = 0$, which is why the familiar form
$\Delta Y = -4\pi R^3 \varepsilon_s/M$ is correct for a bare quark-star surface
and wrong inside a hybrid star. This is the correction of
Ref.~\cite{TakatsyKovacs2020} to Ref.~\cite{Postnikov2010}; using $m$ alone
biases $\Lambda$ for self-bound stars and for hybrid stars with a strong
transition. The implementation returns $\Delta Y = 0$ for
$\Delta\varepsilon \le 0$ or $\hat m \le 10^{-12}$.

\paragraph{How a transition is crossed.} When a Maxwell plateau is present the
integration is split. A second terminal event fires at $\hat p = P_t/P_c$; at
that radius the state $(\hat m, \hat p, Y)$ is taken from the dense output, $Y$
is replaced by $Y + \Delta Y$ from Eq.~\eqref{eq:dY}, the EoS interpolant is
restricted to the low-density branch ($\varepsilon \le 1.01\,\varepsilon_{\rm
low}$), $c_s^2$ is retabulated on that branch, and the integration is resumed
from $r_d$ to the surface. The plateau itself is located beforehand by scanning
the table for a run of points with $|dP|/P$ below tolerance.

\section{The crust}
\label{sec:crust}

A model computes uniform dense matter. Below roughly
$\nB \sim 0.08\;\mathrm{fm^{-3}}$ matter is not uniform, and that regime comes
from a separate tabulated calculation --- here the BPS
crust~\cite{BaymPethickSutherland1971} or a subnuclear slice of a CompOSE
table~\cite{TypelCompOSE2015}. The two must be joined into one
$(P, \varepsilon, \nB)$ table.

\textbf{Splitting on density alone does not work.} The crust and the core are
independent calculations and disagree on $P$ at the same $\nB$: at
$\nB = 0.080\;\mathrm{fm^{-3}}$ BPS gives $P = 0.406\;\mathrm{MeV\,fm^{-3}}$
where one DD2 parametrization gives $0.225$. Concatenating at a chosen
$\nB$ therefore steps \emph{down} in $P$ at the seam. A table whose $P$
decreases with $\nB$ is not an equation of state: $\varepsilon(P)$ becomes
double-valued, the implied $c_s^2$ is negative there, and an integration
crossing it diverges. Three join modes are implemented, and they differ
precisely in how they avoid this.

\begin{description}
\item[\texttt{attach}] Keep crust points with $\nB \le n_t$ AND
  $P < \min P_{\rm core}$, and core points with $\nB > n_t$. The second
  condition is what makes it safe: the join happens where the two $P(\nB)$
  curves \emph{cross}, not at a density fixed in advance, so $P$ cannot step
  down. $n_t$ only bounds how much crust is considered.

\item[\texttt{interpolate}] Blend in $\nB$ over $n_t \pm \delta n$ with
  %
  \begin{equation}
  f(\nB) = \tfrac12\left[1 + \tanh\!\left(\frac{\nB - n_t}{\delta n/2}\right)\right] ,
  \label{eq:blend_n}
  \end{equation}
  %
  applied to $P$ and to the baryon chemical potential
  $\mu_B = (P + \varepsilon)/\nB$, from which
  $\varepsilon = \mu_B \nB - P$ is recovered. Blending $\mu_B$ rather than
  $\varepsilon$ directly is what keeps the blended region thermodynamically
  consistent: the Euler relation is linear in the pair $(P, \mu_B)$ at fixed
  $\nB$, so a convex combination of two consistent states is consistent.

\item[\texttt{maxwell}] Find $P_t$ where $\mu_B^{\rm crust}(P) =
  \mu_B^{\rm core}(P)$ --- the genuine Maxwell condition, equal $\mu_B$ at equal
  $P$ --- and blend in pressure,
  %
  \begin{equation}
  \varepsilon(P) = \tfrac12\left[1 - \tanh\frac{P - P_t}{\delta P}\right]
                     \varepsilon_{\rm crust}(P)
                 + \tfrac12\left[1 + \tanh\frac{P - P_t}{\delta P}\right]
                     \varepsilon_{\rm core}(P) ,
  \label{eq:blend_P}
  \end{equation}
  %
  reducing to a sharp construction as $\delta P \to 0$.
\end{description}

\noindent
Rows with non-finite $P$, $\varepsilon$ or $\nB$, and rows below
$\nB = 10^{-8}\;\mathrm{fm^{-3}}$, are dropped before any of this.

\paragraph{Where the tables come from.} Crust tables are large external data,
neither shipped with the package nor tracked in git. They are resolved by name
--- \texttt{BPS}, \texttt{compose\_sfho\_nYCT},
\texttt{compose\_sfho\_nT0\_beta}, \texttt{compose\_sfho\_nYLS\_trap} --- against
an explicit path, then each directory of \texttt{\$EOS\_CRUST\_DIR} (a search
path, several directories separated by the platform's path separator), then
\texttt{<repo>/data/crust}. A name that cannot be resolved raises with the
file, every directory tried and the variable that adds one. \texttt{have\_crust}
is the predicate a caller uses before falling back: falling back to no crust
moves $M_{\max}$ by about 1\%, so the decision is made visible rather than
silent.

\section{Sequences}
\label{sec:sequences}

A sequence integrates Eq.~\eqref{eq:dmdr}--\eqref{eq:dpdr} once per central
energy density on a grid, normally logarithmic, and returns the columns
$(\varepsilon_c, n_c, P_c, R, M, M_b, k_2, \Lambda)$.

$M_{\max}$ is not read off the grid. The maximum is located on the finite
subset of the sequence --- central states beyond the table's validity integrate
to NaN and must not enter an interpolation --- a window of $\pm 5$ points about
the discrete maximum is fitted with a monotone cubic, and the peak of the fit
is returned together with the index of the last stable point. The stable branch
is everything up to $dM/d\varepsilon_c = 0$; beyond it the configurations are
unstable to radial oscillations and are not stars, so
\texttt{truncate\_to\_stable\_branch} removes them before a mass--radius curve
is drawn.

\section{Uniform rotation}
\label{sec:rotating}

Rotating models are not computed here. They are computed by
RNS~\cite{StergioulasFriedman1995}, a third-party Fortran code driven as a
subprocess; this module writes its input, runs it and parses its output. What
RNS solves is stated first, because a reader cannot check the columns this
module parses without knowing what they are columns of.

\subsection{The formulation RNS solves}
\label{sec:rnsform}

A uniformly rotating, axisymmetric, stationary perfect fluid in the
quasi-isotropic gauge of Komatsu, Eriguchi and
Hachisu~\cite{Komatsu1989}, in the form Cook, Shapiro and
Teukolsky~\cite{CookShapiroTeukolsky1994} put it and RNS implements:
%
\begin{equation}
ds^2 = -e^{\gamma+\rho}\,dt^2 + e^{2\alpha}\big(dr^2 + r^2 d\theta^2\big)
       + e^{\gamma-\rho} r^2\sin^2\!\theta\,\big(d\phi - \omega\,dt\big)^2 ,
\label{eq:kehmetric}
\end{equation}
%
with the four potentials $\gamma$, $\rho$, $\alpha$ and $\omega$ functions of
$(r,\theta)$ alone. The stress tensor is that of a perfect fluid,
$T^{\mu\nu} = (\varepsilon + P)u^\mu u^\nu + P g^{\mu\nu}$, and the rotation is
uniform: the angular velocity $\Omega = u^\phi/u^t$ is one number, not a
profile. $\omega(r,\theta)$ is the frame-dragging rate and is \emph{not} the
same quantity.

Three of the four potentials satisfy elliptic equations which KEH converts to
integral equations with the flat-space Green's functions and solves by
iteration; $\alpha$ then follows by quadrature along a ray. That iteration is
what the accuracy and relaxation parameters of the retry ladder below control.

The matter side is one equation. For a barotrope the Euler equation integrates
to a first integral over the whole star,
%
\begin{equation}
H(r,\theta) + \tfrac12\big(\gamma + \rho\big)
  - \tfrac12\ln\!\big(1 - v^2\big) = \text{const} ,
\qquad
v = \big(\Omega - \omega\big)\,r\sin\theta\;e^{-\rho} ,
\label{eq:bernoulli}
\end{equation}
%
where $v$ is the fluid's proper velocity measured by a zero-angular-momentum
observer and
%
\begin{equation}
H(P) = \int_0^{P}\frac{dP'}{\varepsilon(P') + P'}
\label{eq:enthalpy}
\end{equation}
%
is the specific enthalpy. Eq.~\eqref{eq:enthalpy} is the only place the
equation of state enters, which is why the table handed to RNS carries an
enthalpy column at all, and why this module integrates it on a refined grid
rather than on the table's own rows (Sec.~\ref{sec:rnsinterface}).

A model is specified by two numbers: the central energy density
$\varepsilon_c$ and the axis ratio
%
\begin{equation}
r_{\rm ratio} = \frac{r_p}{r_e} ,
\end{equation}
%
the ratio of the polar to the equatorial \emph{coordinate} radius, with
$r_{\rm ratio} = 1$ the non-rotating star. $\Omega$ is an output, not an input:
the solver finds the rotation rate that makes Eq.~\eqref{eq:bernoulli} hold for
the requested shape. Decreasing $r_{\rm ratio}$ spins the star up until
mass-shedding, where the fluid velocity at the equatorial surface equals the
orbital velocity of a free particle there; that is the Kepler limit, and it
terminates the sequence.

\subsection{Driving the external code}
\label{sec:rnsinterface}

The
binary is located from an explicit argument, then \texttt{\$RNS\_BIN}, then
\texttt{PATH}, then \texttt{<repo>/external/rns}.

Four properties of RNS shape the interface and are worth stating, because each
one is a constraint the caller cannot see:

\begin{itemize}
\item \textbf{The table is capped at 200 rows.} RNS declares
  \texttt{double[201]} and does not check the bound, so a longer table
  overruns silently. A denser table is resampled, not rejected.
\item \textbf{The surface is hardwired} at $\rho = 7.8\;\mathrm{g\,cm^{-3}}$.
  A table that does not reach it is rejected with a message saying to attach a
  crust, rather than being integrated to a wrong answer. The check exists to
  catch a \emph{missing} crust, which is wrong by seven decades, so it is
  applied with a factor-of-ten tolerance.
\item \textbf{The path buffer is \texttt{char[80]}}, filled by
  \texttt{sscanf("\%s")}. The EoS file is therefore symlinked to a short fixed
  name inside the run directory, so the path handed to the solver can never
  approach that limit.
\item \textbf{The enthalpy column must round-trip.} RNS divides it by its own
  value of $c$, $2.9979\times10^{10}\;\mathrm{cm\,s^{-1}}$; writing it with the
  same constant makes the round trip exact instead of accurate to $3\times
  10^{-6}$.
\end{itemize}

\noindent
Non-convergence is retried on a ladder of (accuracy, relaxation factor) pairs,
$(10^{-5}, 1.0)$, $(10^{-6}, 0.8)$, $(10^{-7}, 0.8)$, $(10^{-4}, 0.6)$,
tightening before loosening: most failures here are an iteration oscillating
about the solution, which a smaller step fixes and a larger tolerance only
hides.

The Kepler (mass-shedding) limit and constant-$J$ or constant-$M_0$ sequences
are built by scanning the axis ratio $r_{\rm ratio}$. Along a rotating sequence
the stability boundary is the turning point $\partial M/\partial \varepsilon_c
= 0$ at fixed $J$ or fixed $M_0$~\cite{FriedmanIpserSorkin1988}; the
implementation takes the \emph{first} stationary point, which is the one that
bounds the stable branch.

\section{What a solved star returns}
\label{sec:returns}

\begin{center}
\begin{tabular}{lll}
\toprule
field & quantity & where it comes from \\
\midrule
\texttt{e\_c}   & $\varepsilon_c$ [MeV\,fm$^{-3}$] & the grid point \\
\texttt{n\_c}   & $\nB(P_c)$ [fm$^{-3}$]           & the table \\
\texttt{P\_c}   & $P_c$ [MeV\,fm$^{-3}$]           & $P(\varepsilon_c)$ \\
\texttt{R}      & $R$ [km]                          & $\hat r_s \rsun$ \\
\texttt{M}      & $M$ [$\Msun$]                     & $\hat m(\hat r_s)$ \\
\texttt{M\_b}   & $M_b$ [$\Msun$]                   & Eq.~\eqref{eq:Mb} \\
\texttt{k2}     & $k_2$                             & Eq.~\eqref{eq:k2} \\
\texttt{Lambda} & $\Lambda$                         & Eq.~\eqref{eq:Lambda} \\
\bottomrule
\end{tabular}
\end{center}

\noindent
\texttt{M\_b}, \texttt{k2} and \texttt{Lambda} are \texttt{None} when not
requested; the tidal integration roughly doubles the cost of a point, so it is
optional rather than always on.

A rotating model returns more, and a different object --- \texttt{RotatingResult}
rather than a row of the static sequence. Every field is parsed from RNS output,
so the units are the ones RNS reports rather than the ones used above:

\begin{center}
\begin{tabular}{lll}
\toprule
field & quantity & unit / note \\
\midrule
\texttt{e\_c}      & $\varepsilon_c$                    & MeV\,fm$^{-3}$ \\
\texttt{r\_ratio}  & $r_p/r_e$                          & 1 is non-rotating \\
\texttt{M}         & gravitational mass                  & $\Msun$ \\
\texttt{M\_0}      & rest mass                          & $\Msun$, RNS's own $m_B = 1.66\times10^{-24}$\,g \\
\texttt{R\_e}      & circumferential equatorial radius   & km \\
\texttt{r\_e}, \texttt{r\_p} & coordinate radii         & km; $r_p = r_e\,r_{\rm ratio}$ \\
\texttt{Omega}     & $\Omega$                            & rad\,s$^{-1}$ \\
\texttt{freq}      & spin frequency                      & Hz \\
\texttt{Omega\_K}, \texttt{freq\_K} & mass-shedding rate at the equator \emph{for this model} & rad\,s$^{-1}$, Hz \\
\texttt{J}         & angular momentum                    & dimensionless, $cJ/(G\Msun^2)$ \\
\texttt{I}         & moment of inertia                   & $10^{45}$\,g\,cm$^2$; NaN when non-rotating, where it is undefined \\
\texttt{T\_over\_W} & $T/|W|$                            & rotational over gravitational binding energy \\
\texttt{Phi\_2}    & mass quadrupole moment              & $10^{42}$\,g\,cm$^2$ \\
\texttt{Z\_p}      & polar redshift                      & \\
\texttt{Z\_f}, \texttt{Z\_b} & forward, backward equatorial redshift & \\
\bottomrule
\end{tabular}
\end{center}

\noindent
Two of these are easy to misread. \texttt{Omega\_K} is the mass-shedding rate
\emph{of the model in hand}, not the Kepler frequency of the sequence: a slowly
rotating star reports the $\Omega$ at which \emph{it} would shed, which is not
where the sequence terminates. And \texttt{I} is NaN rather than zero at
$r_{\rm ratio} = 1$, because it is obtained as $J/\Omega$ and both vanish.

A non-converged model carries NaN in every physical field, \texttt{converged =
False} and a \texttt{note} saying why, rather than raising --- CLAUDE.md
section 6's rule that non-convergence is a return value, applied to a sweep that
must survive its individual bad points.

\bibliographystyle{apsrev4-2}
\bibliography{../../../docs/eos}

\end{document}
